\documentclass{article}
\usepackage{geometry}
\usepackage{hyperref}

\title{Course Outline: Economics of Global Shifts}
\author{}
\date{Winter 2025}

\begin{document}

\maketitle

\section*{Course Information}
\begin{itemize}
    \item \textbf{Instructor:} [Muhebullah Karimzada]
    \item \textbf{Course Code:} [ECON 220]
    \item \textbf{Class Schedule:} [Asynchronous]
    \item \textbf{Classroom:} [TBA]
    \item \textbf{Office Hours:} [TBA]
    \item \textbf{Email:} [TBA]
\end{itemize}

\section*{Course Description}
This course explores the concept of global economic shifts and how they impact regions such as Northern British Columbia (NBC). Topics include the effects of technological advancements, demographic changes, trade, environmental factors, and geopolitical dynamics on local and global economies. Using NBC as a case study, students will analyze how regions adapt to and leverage these global shifts for sustainable economic growth. The course will also examine Indigenous perspectives and the role of local communities in navigating economic transitions.

\section*{Course Objectives}
By the end of this course, students will be able to:
\begin{itemize}
    \item Define and understand global economic shifts.
    \item Identify key global drivers such as technology, trade, and demographics.
    \item Analyze the impacts of these shifts on Northern British Columbia's economy.
    \item Evaluate strategies for balancing economic growth with sustainability.
    \item Discuss the role of Indigenous communities in regional development.
\end{itemize}

\section*{Required Texts and Materials}
\begin{itemize}
    \item Stiglitz, J. E. (2002). \textit{Globalization and Its Discontents}.
    \item Bowles, P., \& Wilson, M. (2015). \textit{Resource Communities in a Globalizing Region}.
    \item Prince Rupert Port Authority (2023). \textit{Annual Reports}.
    \item Additional readings will be provided via the course website.
\end{itemize}

\section*{Course Topics and Schedule}

\begin{itemize}
    \item \textbf{Week 1: Introduction to Global Economic Shifts}
    \begin{itemize}
        \item What are global economic shifts?
        \item Key drivers of global shifts (technology, trade, demographics).
    \end{itemize}
    
    \item \textbf{Week 2: The Impact of Technology on Global and Local Economies}
    \begin{itemize}
        \item Technological advancements and automation.
        \item Case studies on technology’s role in reshaping industries in NBC.
    \end{itemize}
    
    \item \textbf{Week 3: Trade and Economic Integration}
    \begin{itemize}
        \item Free trade agreements and their impact on regional economies.
        \item NBC's position within the Asia-Pacific trade network.
    \end{itemize}
    
    \item \textbf{Week 4: Environmental and Climate Changes}
    \begin{itemize}
        \item The global shift towards sustainability.
        \item The role of renewable energy and sustainable practices in NBC's development.
    \end{itemize}
    
    \item \textbf{Week 5: Demographic Changes and their Economic Implications}
    \begin{itemize}
        \item Aging populations and labor force changes.
        \item Outmigration and youth retention in Northern BC.
    \end{itemize}
    
    \item \textbf{Week 6: Geopolitical Dynamics and Economic Shifts}
    \begin{itemize}
        \item Global trade tensions, sanctions, and their effects on NBC.
        \item NBC's reliance on international markets.
    \end{itemize}
    
    \item \textbf{Week 7: Indigenous Perspectives in Economic Development}
    \begin{itemize}
        \item Indigenous community involvement in resource development.
        \item Case study: Joint ventures and Indigenous-led projects.
    \end{itemize}
    
    \item \textbf{Week 8: Case Study: Prince Rupert Port}
    \begin{itemize}
        \item Growth of the port and its economic significance.
        \item Environmental and community impacts.
    \end{itemize}
    
    \item \textbf{Week 9: NBC's Economic Challenges and Opportunities}
    \begin{itemize}
        \item Diversifying the economy beyond resource extraction.
        \item Strategies for sustainable growth and development.
    \end{itemize}
    
    \item \textbf{Week 10: Final Exam Review and Conclusion}
    \begin{itemize}
        \item Recap of key concepts and course discussions.
        \item Preparation for final exam.
    \end{itemize}
\end{itemize}

\section*{Grading and Assessment}
\begin{itemize}
    \item \textbf{Participation:} 10\%
    \item \textbf{Midterm Exam:} 30\%
    \item \textbf{Case Study Presentation:} 20\%
    \item \textbf{Final Exam:} 40\%
\end{itemize}

\section*{Course Policies}
\begin{itemize}
    \item \textbf{Attendance:} Regular attendance is expected. If you are unable to attend a class, please inform the instructor in advance.
    \item \textbf{Late Submissions:} Late assignments will be penalized 5\% per day, unless otherwise arranged with the instructor.
    \item \textbf{Academic Integrity:} Students are expected to adhere to the University’s academic integrity policies. Any form of plagiarism or cheating will not be tolerated.
\end{itemize}

\section*{Discussion Questions for Reflection}
\begin{itemize}
    \item What are the main global trends affecting Northern British Columbia’s economy?
    \item How can Northern British Columbia balance economic growth with environmental sustainability?
    \item What role should Indigenous communities play in the economic development of Northern British Columbia?
\end{itemize}

\end{document}
